\section*{\texorpdfstring{\S\CurrentExerciseGroup}{Section \CurrentExerciseGroup}}
\phantomsection
\addcontentsline{toc}{section}{\texorpdfstring{Exercises for \S\CurrentExerciseGroup}{Exercises for Section \CurrentExerciseGroup}}

\begin{enumerate}
\def\labelenumi{\arabic{enumi})}
\item
  Let \((\lambda_{\alpha})\) be a family of complex point measures on T. For the family \((\lambda_{\alpha})\) to be summable in the vector space \(\mathcal{M}(\mathrm{T};\mathbf{C})\) equipped with the vague topology, it is necessary and sufficient that the family \((|\lambda_{\alpha}|)\) be so.
\item
  Let \((\lambda_{n})\) be a sequence of complex measures on T. For the sequence \((\lambda_{n})\) to be summable in \(\mathcal{M}(\mathrm{T};\mathbf{C})\) equipped with the vague topology, it is necessary and sufficient that, for every function \(f\in \mathcal{K}(\mathrm{T})\) , the series with general term \(\lambda_{n}(f)\) be absolutely convergent (make use of the Banach-Steinhaus theorem; cf.~TVS, III, §4, No.~2, Cor. 2 of Th. 1).
\item
  Let \(\mathbf{T}\) be the compact interval \([0,1]\) of \(\mathbf{R}\) ; for every integer \(n > 0\) , let \(\lambda_{n}\) be the measure
\end{enumerate}

\[
f \mapsto \frac {1}{n} \int_ {0} ^ {1} f (x) \sin 2 n \pi x d x
\]

on T. Show that, in the space \(\mathcal{M}(\mathrm{T})\) equipped with the vague topology (or with the quasi-strong topology (Ch. III, §1, Exer. 8)), the sequence of measures \((\lambda_n)\) is summable, but the sequence \((|\lambda_n|)\) is not. (Observe that if \(\alpha_n = n \cdot \lambda_n(f)\) then \(\sum_{n} \alpha_n^2 < +\infty\) , by making use of Bessel's inequality, and deduce from this that \(\sum_{n} |\lambda_n(f)| < +\infty\) by the Cauchy-Schwarz inequality.)

\begin{enumerate}
\def\labelenumi{\arabic{enumi})}
\setcounter{enumi}{3}
\tightlist
\item
  Let \(\mu\) be a positive measure on T, and \((\nu_{i})_{i\in I}\) a summable family of positive measures on T such that \(\mu \leqslant \sum_{i}\nu_{i}\) . There exists a summable family \((\mu_{i})_{i\in I}\) of measures on T such that \(\mu = \sum_{i}\mu_{i}\) and \(\mu_{i} \leqslant \nu_{i}\) for all \(i \in I\) . (Using Prop. 4, reduce to the case that T is compact, then to the case that \(I = N\) . Then construct the \(\mu_{i}\) by induction.)
\end{enumerate}

\setcounter{exercisegroup}{3}
\edef\CurrentExerciseGroup{\arabic{exercisegroup}}
